% GENERATED FILE. Edit canonical lessons, then run npm run generate:books.
% canonical-source-hash: fnv1a64:dc1429e98cd86624
% canonical-lessons: GE-C110-neben, GE-C110-zwischen, GE-C110-hinter, GE-C110-u-bahn, GE-C110-strassenbahn

\chapter{Next to, Between, Behind, Underground, Tram}
\label{ch:ge-next-to-between-behind-underground-tram}
\hlchaptermodality{110}

\begin{chapteropening}
\textbf{By the end of this chapter:} \emph{I can say next to, between, behind, the underground and a tram.}

Complete the last lesson of chapter 110: I can say next to, between, behind, the underground and a tram.
\end{chapteropening}

\section[neben]{neben — next to}

\label{lesson:GE-C110-neben}



\begin{quote}
Before the new one: say the German for at the back, then the German for opposite.
\end{quote}

\subsection*{You'll want to know: neben}
\textbf{neben} — \textquotedblleft{}next to\textquotedblright{}.

\begin{cousinweb}
\textbf{neben} $\leftarrow$ Middle High German \emph{eneben}, \textquotedblleft{}level with\textquotedblright{}: \emph{in} + the word for \textquotedblleft{}even, level\textquotedblright{}, which EtymDB 2.1 lists beside English \emph{even}. Next to is \textquotedblleft{}on a level with\textquotedblright{}.
\end{cousinweb}

\begin{grammarlens}[title={What you've built}]
One more word for this chapter.
\end{grammarlens}

\subsection*{Your turn}
\begin{itemize}
\raggedright
  \item \emph{Say it:} \emph{neben}
  \item \emph{Say it:} \emph{neben}, once more
  \item \emph{Say it:} say \emph{neben}
  \item \emph{Recall:} say the German for at the back, then the German for opposite, then say \emph{neben} again
\end{itemize}

\begin{tcolorbox}[breakable,colback=teal!4,colframe=teal!35!black,title={Before you move on}]
What does neben mean? (\textquotedblleft{}Next to\textquotedblright{}.) Say it once more.
\end{tcolorbox}
\section[zwischen]{zwischen — between}

\label{lesson:GE-C110-zwischen}



\begin{quote}
Before the new one: say the German for opposite, then the German for next to.
\end{quote}

\subsection*{You'll want to know: zwischen}
\textbf{zwischen} — \textquotedblleft{}between\textquotedblright{}.

\begin{cousinweb}
\textbf{zwischen} $\leftarrow$ Old High German \emph{in zwiskēn}, \textquotedblleft{}among the two\textquotedblright{}, from \emph{zwiski}, \textquotedblleft{}twofold\textquotedblright{}, built on \emph{*twai}, \textquotedblleft{}two\textquotedblright{} (EtymDB 2.1). Between is in the middle of two.
\end{cousinweb}

\begin{grammarlens}[title={What you've built}]
One more word for this chapter.
\end{grammarlens}

\subsection*{Your turn}
\begin{itemize}
\raggedright
  \item \emph{Say it:} \emph{zwischen}
  \item \emph{Say it:} \emph{zwischen}, once more
  \item \emph{Say it:} say \emph{zwischen}
  \item \emph{Recall:} say the German for opposite, then the German for next to, then say \emph{zwischen} again
\end{itemize}

\begin{tcolorbox}[breakable,colback=teal!4,colframe=teal!35!black,title={Before you move on}]
What does zwischen mean? (\textquotedblleft{}Between\textquotedblright{}.) Say it once more.
\end{tcolorbox}
\section[hinter]{hinter — behind}

\label{lesson:GE-C110-hinter}



\begin{quote}
Before the new one: say the German for next to, then the German for between.
\end{quote}

\subsection*{You'll want to know: hinter}
\textbf{hinter} — \textquotedblleft{}behind\textquotedblright{}.

\begin{cousinweb}
\textbf{hinter} $\leftarrow$ Old High German \emph{hintar} $\leftarrow$ Proto-Germanic \emph{*hinder} — inherited, the same word as English \emph{hinder}, \textquotedblleft{}rear\textquotedblright{}, as in \emph{hindquarters}. (EtymDB 2.1)
\end{cousinweb}

\begin{grammarlens}[title={What you've built}]
One more word for this chapter.
\end{grammarlens}

\subsection*{Your turn}
\begin{itemize}
\raggedright
  \item \emph{Say it:} \emph{hinter}
  \item \emph{Say it:} \emph{hinter}, once more
  \item \emph{Say it:} say \emph{hinter}
  \item \emph{Recall:} say the German for next to, then the German for between, then say \emph{hinter} again
\end{itemize}

\begin{tcolorbox}[breakable,colback=teal!4,colframe=teal!35!black,title={Before you move on}]
What does hinter mean? (\textquotedblleft{}Behind\textquotedblright{}.) Say it once more.
\end{tcolorbox}
\section[die U-Bahn]{die U-Bahn — the underground}

\label{lesson:GE-C110-u-bahn}



\begin{quote}
Before the new one: say the German for between, then the German for behind.
\end{quote}

\subsection*{You'll want to know: die U-Bahn}
\textbf{die U-Bahn} — \textquotedblleft{}the underground\textquotedblright{}.

\begin{cousinweb}
\textbf{U-Bahn} is short for \textquotedblleft{}underground railway\textquotedblright{}: the \emph{U} is the first letter of the word for \textquotedblleft{}underground\textquotedblright{}. \emph{Bahn} $\leftarrow$ Old High German \emph{bana} $\leftarrow$ Proto-Germanic \emph{*banō}, \textquotedblleft{}path, cleared way\textquotedblright{} (EtymDB 2.1) — inherited.
\end{cousinweb}

\begin{grammarlens}[title={What you've built}]
One more word for this chapter.
\end{grammarlens}

\subsection*{Your turn}
\begin{itemize}
\raggedright
  \item \emph{Say it:} \emph{die U-Bahn}
  \item \emph{Say it:} \emph{die U-Bahn}, once more
  \item \emph{Say it:} say \emph{die U-Bahn}
  \item \emph{Recall:} say the German for between, then the German for behind, then say \emph{die U-Bahn} again
\end{itemize}

\begin{tcolorbox}[breakable,colback=teal!4,colframe=teal!35!black,title={Before you move on}]
What does die U-Bahn mean? (\textquotedblleft{}The underground\textquotedblright{}.) Say it once more.
\end{tcolorbox}
\section[die Straßenbahn]{die Straßenbahn — a tram}

\label{lesson:GE-C110-strassenbahn}



\begin{quote}
Before the new one: say the German for behind, then the German for the underground.
\end{quote}

\subsection*{You'll want to know: die Straßenbahn}
\textbf{die Straßenbahn} — \textquotedblleft{}a tram\textquotedblright{}.

\begin{cousinweb}
\textbf{Straßenbahn} is the word for \textquotedblleft{}street\textquotedblright{} + \emph{Bahn}, \textquotedblleft{}way, railway\textquotedblright{}: a street railway. The \textquotedblleft{}street\textquotedblright{} word $\leftarrow$ Middle High German \emph{strāze} (EtymDB 2.1), from Latin \emph{strāta}, \textquotedblleft{}paved road\textquotedblright{} — an early loan, like English \emph{street}.
\end{cousinweb}

\begin{grammarlens}[title={What you've built}]
One more word for this chapter.
\end{grammarlens}

\subsection*{Your turn}
\begin{itemize}
\raggedright
  \item \emph{Say it:} \emph{die Straßenbahn}
  \item \emph{Say it:} \emph{die Straßenbahn}, once more
  \item \emph{Say it:} say \emph{die Straßenbahn}
  \item \emph{Recall:} say the German for behind, then the German for the underground, then say \emph{die Straßenbahn} again
\end{itemize}

\begin{tcolorbox}[breakable,colback=teal!4,colframe=teal!35!black,title={Before you move on}]
What does die Straßenbahn mean? (\textquotedblleft{}A tram\textquotedblright{}.) Say it once more.
\end{tcolorbox}
